% Generated by scripts/assemble.py; corrections belong in source/corrections.json.
\examhead{167}{电子科技大学 2024 年数学分析试卷}{158}
% p159-b20; source PDF page 159
\par\Needspace{5\baselineskip}
\blockanchor{p159-b20}
\majorquestion{一}填空题. 每题 5 分, 共 30 分.\par
% p159-b21; source PDF page 159
\blockanchor{p159-b21}
\question{1}计算 $\lim_{x\to 0}\frac{1 - (\cos x)^{\sin x}}{\sqrt{1 + x^3} - 1} =$ \_\_\_\_.\par
% p159-b22; source PDF page 159
\blockanchor{p159-b22}
\question{2}设 $u(x, y, z) = xyz$, 则 $\mathrm{d}^3 u = \underline{\quad}$.\par
% p159-b23; source PDF page 159
\blockanchor{p159-b23}
\question{3}曲线积分 $\oint_{L} (2x^{2} + y^{2} + z) \, \mathrm{d}s =$ \_\_\_\_，其中 $L$ 为 $x^{2} + y^{2} + z^{2} = 1$ 与 $x + y + z = 0$ 的交线.\par
% p159-b24; source PDF page 159
\blockanchor{p159-b24}
\question{4}函数 $z = xy$ 在点 \_\_\_\_ 处的法线与平面 $x + 3y + z + 9 = 0$ 垂直.\par
% p160-b00; source PDF page 160
\blockanchor{p160-b00}
\question{5}级数 $\sum_{n=1}^{\infty} \frac{1}{n \cdot 2^n} =$ \_\_\_\_.\par
% p160-b01; source PDF page 160
\blockanchor{p160-b01}
\question{6}将 $f(x) = \pi - x, x \in (0, \pi)$ 展开成正弦级数 \_\_\_\_, 并写出和函数\par
% p160-b02; source PDF page 160
\par\Needspace{5\baselineskip}
\blockanchor{p160-b02}
\majorquestion{二}计算题. 每题 10 分, 共 60 分.\par
% p160-b03; source PDF page 160
\par\Needspace{6\baselineskip}
\blockanchor{p160-b03}
\question{7}方程组\par
\nopagebreak[4]
% p160-b04; source PDF page 160
\blockanchor{p160-b04}
\displayformula{\left\{ \begin{array}{l} x y + y + y ^ {2} z = 0; \\ x ^ {y} + y z - z ^ {2} + 5 = 0. \end{array} \right.}
% p160-b05; source PDF page 160
\blockanchor{p160-b05}
(可能有误)\par
% p160-b06; source PDF page 160
\blockanchor{p160-b06}
在点 $P(1, -2, 1)$ 附近能否唯一确定隐函数组？\par
% p160-b07; source PDF page 160
\par\Needspace{6\baselineskip}
\blockanchor{p160-b07}
\question{8}计算曲线积分\par
\nopagebreak[4]
% p160-b08; source PDF page 160
\blockanchor{p160-b08}
\displayformula{\int_ {L} \frac {(x - y) \mathrm{d} x + (x + 4 y) \mathrm{d} y}{x ^ {2} + 4 y ^ {2}}.}
% p160-b09; source PDF page 160
\blockanchor{p160-b09}
其中 $L: x^{2} + y^{2} = 1$, 取逆时针方向.\par
% p160-b10; source PDF page 160
\blockanchor{p160-b10}
\question{9}求 $x + 2y = 1$ 与 $x^{2} + 2y^{2} + z^{2} = 1$ 相交曲线上到原点距离最小的点.\par
% p160-b11; source PDF page 160
\blockanchor{p160-b11}
\question{10}判断数项级数 $\sum_{n=1}^{\infty} (-1)^{n-1} \frac{\sin^2 n}{n}$ 的敛散性, 若收敛需要判断是条件收敛还是绝对收敛.\par
% p160-b12; source PDF page 160
\blockanchor{p160-b12}
\question{11}计算积分 $I(y) = \int_{0}^{+\infty} e^{-x^2} \cos 2xy \, \mathrm{d}x$，其中 $y \in \mathbb{R}$.\par
% p160-b13; source PDF page 160
\blockanchor{p160-b13}
\question{12}计算二重积分 $\iint_{D} (\sqrt{x^2 + y^2} + y) \, \mathrm{d}x \, \mathrm{d}y$, 其中 $D: x^2 + y^2 = 4$ 与 $(x + 1)^2 + y^2 = 1$ 所围成的区域.\par
% p160-b14; source PDF page 160
\par\Needspace{5\baselineskip}
\blockanchor{p160-b14}
\majorquestion{三}证明题. 每题10分, 共30分.\par
% p160-b15; source PDF page 160
\blockanchor{p160-b15}
\question{13}已知数列  $\{a_{n}\}$  满足  $a_{n}=\sin1+\frac{\sin2}{2!}+\frac{\sin3}{3!}+\cdots+\frac{\sin n}{n!}$ .\par
% p160-b16; source PDF page 160
\blockanchor{p160-b16}
(1) 证明: $\{a_{n}\}$ 有界, 但是 $\{a_{n}\}$ 不单调.\par
% p160-b17; source PDF page 160
\blockanchor{p160-b17}
(2) 证明: 数列 $\{a_{n}\}$ 收敛.\par
% p160-b18; source PDF page 160
\blockanchor{p160-b18}
\question{14}证明: 函数项级数 $\sum_{n=1}^{\infty} \frac{x^{2}}{(1 + x^{2})^{n}}$ 在 $(- \infty, +\infty)$ 上非一致收敛.\par
% p160-b19; source PDF page 160
\blockanchor{p160-b19}
\question{15}证明: 反常积分 $\int_{0}^{+\infty} \frac{x}{1 + x^{5} \sin^{2} x} \mathrm{d}x$ 收敛.\par
% p160-b20; source PDF page 160
\par\Needspace{5\baselineskip}
\blockanchor{p160-b20}
\majorquestion{四}综合题. 每题 15 分, 共 30 分.\par
% p160-b21; source PDF page 160
\par\Needspace{6\baselineskip}
\blockanchor{p160-b21}
\question{16}已知 $\lambda > 0, \beta \in (0,1), \{a_n\}$ 为正数列，且满足\par
\nopagebreak[4]
% p160-b22; source PDF page 160
\blockanchor{p160-b22}
\displayformula{\liminf _ {n \to \infty} n ^ {\beta} \left(\frac {a _ {n}}{a _ {n + 1}} - 1\right) = \lambda .}
% p160-b23; source PDF page 160
\blockanchor{p160-b23}
证明: 对任意的 k > 0, 有  $\lim_{n \to \infty} n^{k} a_{n} = 0$ .\par
% p160-b24; source PDF page 160
\par\Needspace{6\baselineskip}
\blockanchor{p160-b24}
\question{17}已知函数 $f(x) \in C^2[0,1]$，且满足 $f(0)f(1) \geq 0$. 证明：\par
\nopagebreak[4]
% p160-b25; source PDF page 160
\blockanchor{p160-b25}
\displayformula{\int_ {0} ^ {1} | f ^ {\prime} (x) | \mathrm{d} x \leq 2 \int_ {0} ^ {1} | f (x) | \mathrm{d} x + \int_ {0} ^ {1} | f ^ {\prime \prime} (x) | \mathrm{d} x.}
